% Раздзел 1 — метадалогія devops
\section{Метадалогія Devops}

\subsection{Паняцце Devops}

Паняцце Devops не мае агульнапрынятага азначэння.
Азначэнне Devops змяняецца ў залежнасці ад часу, метадаў укаранення,
пастаўленых мэт унутры кампаніі.
У табліцы \ref{table:Definition of Devops} прыведзены некаторыя
азначэнні, каторыя сустракаюцца ў літаратуры.

\begin{table}[htbp]
    \caption{Варыянты азначэнняў Devops}
    \begin{tabular}{|p{0.70\textwidth}|p{0.24\textwidth}|} 
        \hline
        \textbf{Азначэнне Devops}
        &
        \textbf{Сэнс азначэння}  \\ 
        \hline
        Devops -- гэта набор практык, які прызначыны для
        памяншэння часу паміж прыняццямі змен ў сістэме і
        ўнясення іх ў працэс вытворчасці з гарантыяй
        высокай якасці[\ref{book:DevOps: A Software Architect's Perspective}]
        &
        Арыентацыя на дасягненне мэты (хуткая дастаўка
        якаснага праграмнага забеспячэння) \\ 
        \hline
        Devops -- гэта культурны рух, які змяняе адносіны людзей
        да працы і да яе вынікаў[\ref{book:Effective Devops}]
        & Арыентацыя на супрацоўніцтва \\ 
        \hline
        Devops -- камбінацыя культуры, практык і інструментаў,
        каторая павялічвае здольнасць арганізацыі
        пастаўляць праграмы і сервісы
        з высокай хуткасцю: развіццё і паляпшэнне прадукцыі
        ў больш хуткім тэмпе,
        чым арганізацыі, якія карыстаюцца традыцыйным спосабам распрацоўкі
        праграмнага забеспячэння і
        кіравання працэсамі інфраструктуры.[\ref{site:aws.amazon.com/devops}]
        & Арыентацыя на супрацоўніцтва і дасягненне мэты
        (паскарэнне тэмпаў развіцця прадукту) \\
        \hline
        Devops -- набор практык, накіраваны на актыўнае ўзаемадзеянне
        спецыялістаў па распрацоўцы і
        спецыялістаў па інфармацыйна-тэхналагічнаму
        абслугоўванню і ўзаемную інтэграцыю
        іх працоўных працэсаў
        адно ў другое[\ref{site:ru.wikipedia.org/wiki/devops}]
        &
        Арыентацыя на супрацоўніцтва \\
        \hline
    \end{tabular}
    \label{table:Definition of Devops}
\end{table}

З табліцы \ref{table:Definition of Devops} можам заключыць,
што паняцце Devops уключае ў сябе як культурны, так і
тэхнічны складнікі.
Пры гэтым набор інструментаў, такіх як Chef альбо Docker,
з'яўляецца складнікам Devops дзякуючы спосабам іх прымянення,
а не ў сілу характарыстык саміх інструментаў.[\ref{book:Effective Devops}]
Такім чынам неабходна адзначыць, што ў першую чаргу Devops разглядаецца
як змяненне культуры арганізацыі,
а іменна адносін паміж камандамі (у прыватнасці каманды распрацоўшчыкаў і
каманды тэхнічнага абслугоўвання).

\subsection{Мадэль CALMS}

Нягледзячы на разнастайнае тлумачэнне тэрміна, Devops мае
пэўныя агульнапрызнаныя элементы.
Гэтыя элементы складаюць мадэль CALMS[\ref{article:Challenges in Adopting a Devops}] (малюнак \ref{fig:CALMS model}):
\begin{enumerate}
    \item Culture (культура),
    \item Automation (аўтаматызацыя),
    \item Lean (ашчаднасць),
    \item Measurement (мера),
    \item Sharing (абмен).
\end{enumerate}

\begin{figure}[h!]
    \includegraphics[width=\linewidth]{CALMS.pdf}
    \caption{Мадэль CALMS}
    \label{fig:CALMS model}
\end{figure}

\vspace{-\baselineskip}
\subsubsection{Kyльтура.}

Джэніфэр Дэвіс і Кэтрын Дэніэлс у сваёй кнізе[\ref{book:Effective Devops}]
падкрэсліваюць, што не трэба ігнараваць нормы і каштоўнасці
культур і міжасобасных узаемадзеянняў, звяртаючы ўвагу толькі
на карыстанне інструментамі.

Яны сцвярджаюць, што без развязвання міжасобасных і
міжгрупавых канфліктаў, якія ўзнікаюць у арганізацыі,
немагчыма пабудаваць трывалыя адносіны, якія і прыводзяць
да фар\-мі\-ра\-ван\-ня Devops-асяроддзя.

Джэніфэр Дэвіс і Кэтрын Дэніэлс падзяляюць культуры на 2 пункты:

\begin{enumerate}
    \item Супрацоўніцтва
    \item Блізкасць
\end{enumerate}

Супрацоўніцтва -- гэта працэс дасягнення пастаўленай мэты
пры дапамозе ўзаемадзеяння паміж некалькімі людзьмі.
Галоўным прынцыпам руху Devops стаў кааперацыя паміж
камандамі па распрацоўцы і эксплуатацыі праграмнага забеспячэння.
Пры гэтым перад тым, як арганізаваць паспяховае
ўзаемадзеянне паміж камандамі, каторыя маюць розныя мэты,
неабходна наладзіць групавую работу ў адной камандзе.
Так, калі ў камандах не ўпарадкавана работа на індывідуальным
і групавым узроўнях, немагчыма дасягнуць паспяховага ўзаемадзеяння
паміж камандамі.

Апроч развіцця і падтрымкі адносін супрацоўніцтва паміж
асобнымі людзьмі, групамі і аддзеламі ўнутры арганізацыі,
важны трывалыя ўзаемаадносіны ў сферы дзеянасці.

Блізкасць -- гэта працэс фарміравання ўзаемаадносін паміж
камандамі, які садзейнічае свабодзе выбару розных мэт альбо
паказальнікаў пры захаванні агульных мэт арганізацыі.
Пры гэтым аблягчаецца развіццё эмпаціі і навучання ў розных групах.
Блізкасць можа праяўляцца на ўзроўні ўзаемаадносін паміж
арганізацыямі, што дазваляе ім дзяліцца гісторыямі і вучыць адно
аднаго. У выніку ствараецца калектыўная база культурных і тэхнічных ведаў.
%Уцягнуць персанал тэхнічнай эксплуатацыі ў працэсы
%распрацоўкі праграм. Акрамя таго, забяспечыць іх уклад на
%сходзе для планавання, агляду праектаў, каб яны мелі магчымасць
%выказаць уласныя ідэі ўжо на ранніх стадыях працэсу
%распрацоўкі.
%Humble і Molesky (2011 год) заўважылі, што ўзаемадзеянне з ка\-ман\-дамі
%тэхнічнай эксплуатацыі неабходна для каманд распрацоўшчыкаў.
%Абодва бакі павінны ў аднолькавай меры ўдзельнічаць у аналізе знаходжання
%першапрычын ў выпадку вытворчых інцыдэнтаў, у знаходжанні спосабаў
%развязвання праблемы як на этапах распрацоўкі, так і дастаўкі прадукту кліенту.

%Без разумення працэсаў, каторыя адбываюцца ўнутры іншых каманд
%(тэхнічнай эксплуатацыі, тэсціравання), немагчыма дабіцца
%%палітыкі неабвінавачвання ў кампаніі.
%Палітыка неабвінавачвання прадугледжвае прызнанне чалавечых памылак і
%разглядае іх у якасці магчымасці палепшыць прадукцыйнасць кампаніі
%пры дапамозе аналізу першапрычын і наступнага недапушчэння
%аналагічных памылак.
%Калі супрацоўнікі кампаніі не баяцца быць абвінавачанымі і
%пакаранымі за памылкі, адкрываецца шлях да новых інавацый,
%павялічваецца матывацыя і адданасць кампаніі.
%
%У асяроддзі супрацоўніцтва знікае такая праблема, як перакладванне
%адказнасці на іншыя каманды, дазваляе на ранніх стадыях знаходзіць
%праблемы ў праграмных прадуктах або паслугах.

\subsubsection{Аўтаматызацыя.}

Для таго, каб пабудаваць новы тып культуры ў арганізацыі,
неабходна змяняць працэсы.
Згодна з Гамбэлам і Малескім[\ref{book:Lean Enterprise}],
гэта становіцца магчымым пры дапамозе Devops інструментаў і 
інфраструктуры як код. Мэтай практыкі аўтаматызацыі ў Devops
ёсць дасягнення памяншэння часу выпуску і паскарэння зваротнай
сувязі на змяненні.
Гэта азначае выкарыстоўванне працэсу разгортвання,
які ахоплівае ўсе змяненні, якія могуць быць унесены любой з каманд.
Працэс разгортвання ўключае ў сябе працэсы распрацоўкі, зборкі,
тэсціравання і разгортвання сістэмы.

Аўтаматызацыя працэсаў дазваляе:
\begin{enumerate}
    \item паскорыць распрацоўку праграмнага забеспячэння;
    \item паменшыць імавернасць чалавечай памылкі;
    \item паскорыць атрыманне зваротнай сувязі пры ўнясенні змяненняў
          на этапах распрацоўкі праграмнага забеспячэння.
\end{enumerate}

\subsubsection{Ашчаднасць.}
Ашчаднае мысленне з'яўляецца адной з ключавых частак Devops.

Марыя і Том Попендык[\ref{book:Lean Software Development}]
сфармулявалі рад прынцыпаў для выканання ашчаднага мыс\-лен\-ня ў
распрацоўкі праграмнага забеспячэння.

Першы прынцып -- памяншэнне адходаў.
Адходы -- гэта ўсё, што не дабаўляе цэннасці канчатковаму прадукту.
А цэннасць прадукту заўсёды разглядаецца з пункту гледжання кліента.
Адходамі можа быць лішняя дакументацыя, дадатковыя праграмныя магчымасці,
каторыя не былі яўна запатрабаваныя кліентам,
перадача распрацоўкі ад адной групы іншай.
Галоўная мэта -- дастаўка акурат таго, што хацелі спажыўцы прадукту.
Перавышэнне часу распрацоўкі таксама разглядаецца ў якасці адходаў.

Другі і трэці прынцыпы -- гэта узмацненне навучання і прыняцце рашэнняў як
мага пазней. У кантэксце распрацоўкі праграмнага забеспячэння гэта
значыць, што ітэрацыйныя цыклы распрацоўкі і прыняцце канчатковых
рашэнняў развіцця прадукту адбываецца толькі тады, калі гэта неабходна.
Навучанне ўзмацняецца пры дапамозе пастаянных ацэнак і паляпшэнняў працэсаў.
Прыняцце рашэнняў як мага пазней становіцца эфектыўным у сітуацыях,
звязаных з павелічэннем узроўню нявызначанасці.
Ва ўмовах дынамічнага развіцця рынку, лепшымі рашэннямі з'яўляюцца тыя
рашэнні, якія былі зроблены на падставе фактаў, а не здагадак.
У межах праграмнага забеспячэння гэта азначае, што сістэмы павінны
ўключаць магчымасць унясення змяненняў.


Чацвёрты прынцып -- найхутчэйшая дастаўка.
Хуткая дастаўка дазваляе паскорыць цыкл зваротнай сувязі, што
дазваляе своечасова вызначыць перавагі і недахопы ў новых
версіях прадукту.
Такім чынам хуткая дастаўка становіцца адным з галоўных кампанентаў
метадалогіі Devops.


Пяты і шосты прынцыпы -- упаўнаважванне каманд і пабудова цэласнасці.
Тэхнічныя рашэнні павіны прынімацца людзьмі, якія будуць несці
адказнасць за выкананне планаў, так як пры рашэннях адкладзеных
да апошняй хвіліны, не існуе часу для арганізацыі дзейнасці
працаўнікамі вышэйшах уладай.
Цэласнасць праграмнага забеспячэння стварае надзейную архітэктуру,
высокую практычнасць, рамонтапрыдатнасць, адаптыўнасць
і пашыральнасць[\ref{book:Lean Software Development}].

Апошні прынцып разглядае цэласнасць усёй сістэмы працы.
Цэласнасць у складаных сістэмах патрабуе глыбокага
даследвання ў шматлікіх сферах.
З пункту гледжання Devops гэта азначае, што каманды
распрацоўшчыкаў і тэхнічнай эксплуатацыі павінны
узаемадзейнічаць на працягу ўсяго цыклу распрацоўкі
праграмнага забеспячэння, а не аб'ядноўваць вынікі
працы кожнай каманды ў канцы цыкла.

Прынцыпы ашчаднай распрацоўкі праграмнага забеспячэння
дазволілі перайсці да больш гібкіх мадэлей распрацоўкі
праграмнага забеспячэння.
Метадалогіі на прынцыпах Agile распрацоўкі прывялі да
больш хуткага рэагавання на новыя патрабаванні ў параўнанні
з традыцыйнымі мадэлямі распрацоўкі.

\subsubsection{Мера}
Гамбэл і Малескі[\ref{book:Lean Enterprise}] вызначаюць меру
ў Devops, як высокаўзроўневы кантроль бізнес-метрык,
напрыклад даход альбо колькасць аперацый за адзінку часу.
На больш нізкім узроўні гэта патрабуе дбайнага выбару ключавых
паказчыкаў ў працэсе дастаўкі: час працэсу дастаўкі,
уплыў новых рэлізаў на стабільнасць сістэмы.

Мера ўплывае на тое, як людзі працуюць.
Візуалізацыя метрык дазваляе дакладна
вызначыць прадукцыйнасць каманды ў любы момант часу.
Метрыкі таксама дапамагаюць вызначыць
вузкія месцы ў працэсах распрацоўкі, што ў сваю чаргу
дапамагае падтрымліваць ашчаднасць распрацоўкі.

У метадалогіі Devops могуць выкарыстоўвацца
такія метрыкі як[\ref{site:stackify.com/15-metrics}]:
\begin{enumerate}
    \item частата разгортвання праграмнага забеспячэння
          (як часта адбываецца разгортванне
           новых рэлізаў праграмнага забеспячэння);
    \item аб'ём змяненняў (як шмат змяненняў, выпраўленняў памылак
          уваходзіць у новы рэліз);
    \item час разгортвання (час, які неабходны, для
          разгортвання новага рэлізу);
    \item час працэсу распрацоўкі (час ад пачатку
          стварэння прадукту да яго разгортвання);
    \item працэнт праходжання аўтаматычных тэстаў
          (дазваляе вызначыць як часта змяненні кода
           прыводзяць да памылак у праграмным забеспячэнні);
    \item каэфіцыент памылак (вызначае колькасць памылак,
          якія не былі вызначаны падчас тэсціравання
          праграмнага забеспячэння перад разгортваннем);
    \item сярэдні час знаходжання памылак (як хутка
          знаходзяцца памылкі пасля разгортвання);
    \item сярэдні час аднаўлення сістэмы (як хутка
          выпраўляюцца памылкі пасля іх знаходжання).
\end{enumerate}
%
%З вышэй прыведзеных азначэнняў тэрміна Devops можам вызначыць,
%што, хаця не існуе адзінага разумення метадалогіі Devops,
%у кожным выпадку значная ўвага надаецца культуры ўнутры арганізацыі,
%узаемадзеянню спецыялістаў паміж сабой.
%Дадзеная асаблівасць тлума\-чыц\-ца тым, што метадалогія Devops сфармавалася
%на аснове гібкіх метадалогій, у якіх адбываецца пераход фокуса на
%асобных людзей, узаемадзеянне і супрацоўніцтва.
%
%%Аднак рух Devops пашырае прынцыпы гібкай распрацоўкі праграм
%і прымяняе іх на ўзроўні арганізацыі ў цэлым,
%у той час як іншыя гібкія метадалогіі акцэнтавалі ўвагу толькі на
%распрацоўшчыках праграм.
%
%Нягледзячы на разнастайнае тлумачэнне тэрміна, Devops мае пэўныя
%культурныя прын\-цы\-пы. Гэтыя прынцыпы складаюць мадэль CALMS:
%Culture (культура), Automation (аўтаматызацыя),
%Lean (беражлівасць), Measurement (мера), Sharing (абмен).%
%\ref{article_challenges_in_adopting_a_Devops}
%
%Мадэль CALMS прадстаўлена на малюнку \ref{fig:CALMS model}.
%
%\begin{figure}[h!]
%    \includegraphics[width=\linewidth]{CALMS.pdf}
%    \caption{мадэль CALMS}
%    \label{fig:CALMS model}
%\end{figure}
%
%\vspace{-\baselineskip}
%\subsubsection{Мадэль CALMS -- культура.}
%
%Прыняцце метадалогіі Devops пачалося з культуры.
%У той час як матэрыяльная частка Devops змяшчае новыя працэсы і
%інструменты для бесперапыннай дастаўкі і інтэграцыі,
%Devops, усяго толькі, папулярнае слова,
%калі "культура арганізацыі" не ахоплівае ўсе рэчы,
%што ідуць у камплекце (Wilsenach 2015).
%
%Галоўная ідэя Devops -- змяніць культуру арганізацыі з
%закрытых скрынь (асобныя ка\-ман\-ды) на адкрыты спосаб супрацоўніцтва
%паміж імі.
%Уцягнуць персанал тэхнічнай эксплуатацыі ў працэсы
%распрацоўкі праграм. Акрамя таго, неабходна іх прысутнасць на
%сходзе для планавання, агляду праектных каманд, каб яны маглі
%падзяліцца ўласнымі ідэямі і ведамі ўжо на ранніх стадыях працэсу
%распрацоўкі.
%Humble і Molesky (2011 год) заўважылі, што абмен з ка\-ман\-дамі
%тэхнічнай эксплуатацыі неабходны для распрацоўшчыкаў, і яны
%павінны ў аднолькавай меры ўдзельнічаць у аналізе знаходжання
%першапрычын ў выпадку вытворчых інцыдэнтаў.
%
%Без разумення працэсаў, каторыя адбываюцца ўнутры іншых каманд
%(тэхнічнай эксплуатацыі, тэсціравання), немагчыма дабіцца
%палітыкі неабвінавачвання ў кампаніі.
%Палітыка неабвінавачвання прадугледжвае прызнанне чалавечых памылак і
%разглядае іх у якасці магчымасці палепшыць прадукцыйнасць кампаніі
%%пры дапамозе аналізу першапрычын і наступнага недапушчэння
%аналагічных памылак.
%Калі супрацоўнікі кампаніі не баяцца быць абвінавачанымі і
%пакаранымі за памылкі, адкрываецца шлях да новых інавацый,
%павялічваецца матывацыя і адданасць кампаніі.
%
%У асяроддзі супрацоўніцтва знікае такая праблема, як перакладванне
%адказнасці на іншыя каманды, дазваляе на ранніх стадыях знаходзіць
%праблемы ў праграмных прадуктах або паслугах.
%
%\subsubsection{Мадэль CALMS -- аўтаматызацыя.}
